\chapter{ทฤษฎีและงานวิจัยที่เกี่ยวข้อง}

\section{เอนไซม์เป้าหมายและค่าฤทธิ์ทางชีวภาพ}

\subsection{SARS-CoV-2 Main Protease}

เอนไซม์ SARS-CoV-2 Main Protease (\Mpro) เป็นเอนไซม์ชนิดซิสเตอีนโปรตีเอสที่เข้ารหัสอยู่ในบริเวณ nsp5 ของ Replicase Polyprotein 1ab เอนไซม์นี้ทำงานในรูปไดเมอร์ (Dimer) และมีตำแหน่งเร่งปฏิกิริยา (Catalytic Dyad) ประกอบด้วยกรดอะมิโนฮิสทิดีนที่ตำแหน่ง 41 และซิสเตอีนที่ตำแหน่ง 145

หน้าที่หลักของ \Mpro\ คือการตัดสายพอลิโปรตีน pp1a และ pp1ab ที่ไวรัสสร้างขึ้นให้กลายเป็นโปรตีนย่อยที่ทำงานได้ โดย \Mpro\ รับผิดชอบการตัดอย่างน้อย 11 ตำแหน่ง การยับยั้งเอนไซม์นี้จึงหยุดกระบวนการเพิ่มจำนวนของไวรัสได้โดยตรง Jin และคณะ (2020) ได้รายงานโครงสร้างผลึกของ \Mpro\ พร้อมการค้นพบสารยับยั้ง ซึ่งเป็นงานพื้นฐานสำคัญที่กระตุ้นการวิจัยสารยับยั้ง \Mpro\ อย่างกว้างขวาง

ความสำคัญเชิงเภสัชวิทยาของ \Mpro\ อยู่ที่การที่มนุษย์ไม่มีเอนไซม์ที่มีความคล้ายคลึงเชิงโครงสร้างและหน้าที่ ทำให้สารยับยั้งที่ออกแบบมาจำเพาะต่อ \Mpro\ มีโอกาสเกิดพิษต่อเซลล์เจ้าบ้านต่ำ ตัวอย่างยาที่พัฒนาสำเร็จบนหลักการนี้คือ Nirmatrelvir ซึ่งเป็นองค์ประกอบหลักของยา Paxlovid

\subsection{ค่า \ICfifty\ และการแปลงเป็น \pICfifty}

ค่า \ICfifty\ (Half Maximal Inhibitory Concentration) เป็นตัวชี้วัดความเข้มข้นของสารที่ใช้ยับยั้งการทำงานของเอนไซม์ได้ร้อยละ 50 ข้อมูลในฐานข้อมูล ChEMBL ถูกจัดเก็บในหน่วยนาโนโมลาร์ (nM) และมีการกระจายตัวแบบเบ้ขวาอย่างรุนแรง (Right-skewed) เนื่องจากค่าครอบคลุมหลายลำดับขนาด

การแปลงเป็นค่า \pICfifty\ ช่วยลดความเบ้ของข้อมูลและทำให้ตัวแปรตามมีการกระจายตัวใกล้เคียงปกติ ซึ่งเหมาะสมกับการสร้างแบบจำลองการถดถอย โดยมีความสัมพันธ์ทางคณิตศาสตร์ดังสมการที่ \ref{eq:pic50}

\begin{equation}
\pICfifty = -\log_{10}\left(\ICfifty[\mathrm{M}]\right) = 9 - \log_{10}\left(\ICfifty[\mathrm{nM}]\right)
\label{eq:pic50}
\end{equation}

ตัวอย่างเช่น สารที่มี \ICfifty\ = 100 nM จะมี \pICfifty\ $= 9 - \log_{10}(100) = 9 - 2 = 7.0$ ค่า \pICfifty\ ที่สูงกว่าหมายถึงสารมีฤทธิ์แรงกว่า

\begin{warnnote}[ข้อควรระวังเรื่องหน่วย]
ในบางแหล่งข้อมูลค่า \pICfifty\ ถูกคำนวณบนฐานหน่วยไมโครโมลาร์ ($\mu$M) ซึ่งจะมีความสัมพันธ์ $\pICfifty(\mu\mathrm{M\ basis}) = \pICfifty(\text{มาตรฐาน}) - 6$ การนำข้อมูลจากต่างแหล่งมารวมกันจึงต้องตรวจสอบหน่วยฐานให้ตรงกันก่อนเสมอ
\end{warnnote}

\section{การแทนรหัสโครงสร้างเคมีและคุณลักษณะโมเลกุล}

\subsection{SMILES}

SMILES (Simplified Molecular Input Line Entry System) เป็นระบบแทนโครงสร้างทางเคมีด้วยสายอักขระ ทำให้สามารถจัดเก็บและประมวลผลโครงสร้างโมเลกุลด้วยคอมพิวเตอร์ได้อย่างมีประสิทธิภาพ โมเลกุลเดียวกันอาจเขียนเป็น SMILES ได้หลายรูปแบบ จึงจำเป็นต้องแปลงให้อยู่ในรูปมาตรฐาน (Canonical SMILES) ก่อนการเปรียบเทียบหรือยุบข้อมูลซ้ำ

\subsection{Morgan Fingerprints (ECFP4)}

Morgan Fingerprints เป็นอัลกอริทึมเข้ารหัสโครงสร้าง 2 มิติที่พัฒนาโดย Rogers และ Hahn (2010) หลักการทำงานคือการพิจารณาสภาพแวดล้อมรอบอะตอมแต่ละตัวภายในรัศมีที่กำหนด แล้วแปลงเป็นเวกเตอร์บิตความยาวคงที่

เมื่อกำหนดรัศมี $r = 2$ อัลกอริทึมจะพิจารณาอะตอมข้างเคียงที่อยู่ห่างออกไปไม่เกิน 2 พันธะ ซึ่งเทียบเท่ากับเส้นผ่านศูนย์กลาง 4 พันธะ จึงเรียกว่า ECFP4 (Extended-Connectivity Fingerprint, diameter 4) ผลลัพธ์คือเวกเตอร์บิตขนาด 1{,}024 บิต โดยแต่ละบิตมีค่าเป็น 0 หรือ 1 บ่งบอกการมีอยู่หรือไม่มีของโครงสร้างย่อยแบบใดแบบหนึ่ง

ข้อได้เปรียบสำคัญของ Morgan Fingerprints คือความสามารถในการจับคู่ย้อนกลับ (Bit Mapping) กล่าวคือ เมื่อทราบว่าบิตใดมีความสำคัญ สามารถย้อนกลับไปวาดภาพโครงสร้างย่อยที่บิตนั้นเป็นตัวแทนได้ ทำให้ผลการวิเคราะห์สามารถตีความเชิงเคมีได้

\begin{warnnote}[ข้อจำกัด]
Morgan Fingerprints ในการตั้งค่าปกติไม่บันทึกข้อมูลสเตอริโอเคมี (Stereochemistry) ทำให้สเตอริโอไอโซเมอร์ที่มีโครงสร้างเชิงเส้นเหมือนกันแต่มีการจัดเรียงในสามมิติต่างกันได้เวกเตอร์บิตที่เหมือนกันทุกประการ ข้อจำกัดนี้จะปรากฏผลจริงในบทที่ 5
\end{warnnote}

\subsection{กฎของลิพินสกีและ Lipinski Descriptors}

Lipinski และคณะ (2001) เสนอกฎห้าประการ (Rule of Five) สำหรับประเมินความเป็นไปได้ที่สารจะพัฒนาเป็นยารับประทาน โดยสารที่มีคุณสมบัติเป็นยาที่ดีมักมีลักษณะดังตารางที่ \ref{tab:lipinski}

\begin{table}[H]
\centering
\caption{กฎห้าประการของลิพินสกีและความหมายของแต่ละ Descriptor}
\label{tab:lipinski}
\begin{tabular}{L{4.2cm} C{2.6cm} L{6.4cm}}
\toprule
\textbf{Descriptor} & \textbf{เกณฑ์} & \textbf{ความหมาย}\\
\midrule
MW (Molecular Weight) & $\leq 500$ Da & น้ำหนักโมเลกุล\\
LogP & $\leq 5$ & ค่าสัมประสิทธิ์การแบ่งน้ำ--ออกทานอล บ่งบอกความชอบไขมัน\\
HBD (H-Bond Donors) & $\leq 5$ & จำนวนอะตอมผู้ให้พันธะไฮโดรเจน\\
HBA (H-Bond Acceptors) & $\leq 10$ & จำนวนอะตอมผู้รับพันธะไฮโดรเจน\\
\bottomrule
\end{tabular}
\end{table}

คุณลักษณะเหล่านี้เป็นค่าตัวเลขต่อเนื่องที่มีสเกลแตกต่างกันมาก เช่น MW อยู่ในหลักร้อย ขณะที่ HBD อยู่ในหลักหน่วย จึงจำเป็นต้องปรับมาตรฐานก่อนใช้กับแบบจำลองที่ไวต่อสเกล เช่น SVR

\subsection{คุณลักษณะแบบผสม (Hybrid Features)}

การรวม Morgan Fingerprints ซึ่งเป็นข้อมูลไบนารีเชิงโครงสร้าง เข้ากับ Lipinski Descriptors ซึ่งเป็นข้อมูลตัวเลขต่อเนื่องเชิงคุณสมบัติ ทำให้ได้ชุดคุณลักษณะที่ครอบคลุมทั้งมิติโครงสร้างเฉพาะและมิติคุณสมบัติภาพรวมของโมเลกุล รวมเป็น $1{,}024 + 4 = 1{,}028$ คุณลักษณะ

ความท้าทายของการรวมข้อมูลสองประเภทนี้คือ ข้อมูลไบนารีไม่ควรถูกปรับมาตรฐานเพราะมีค่า 0/1 อยู่แล้ว ขณะที่ข้อมูลตัวเลขต่อเนื่องจำเป็นต้องปรับ จึงต้องใช้กลไก \texttt{ColumnTransformer} เพื่อจัดการคอลัมน์แต่ละกลุ่มแยกจากกัน

\section{อัลกอริทึมการเรียนรู้ของเครื่อง}

\subsection{Dummy Regressor}

แบบจำลองพื้นฐานที่ทำนายผลลัพธ์เป็นค่าเฉลี่ยของชุดข้อมูลฝึกเสมอ โดยไม่พิจารณาคุณลักษณะใด ๆ ใช้เป็นเกณฑ์มาตรฐานเปรียบเทียบขั้นต่ำ (Baseline Benchmark) หากแบบจำลองที่ซับซ้อนกว่าไม่สามารถทำได้ดีกว่า Dummy Regressor แสดงว่าคุณลักษณะที่ใช้ไม่มีสารสนเทศเพียงพอ หรือกระบวนการสร้างแบบจำลองมีข้อผิดพลาด ในทางทฤษฎี Dummy Regressor จะให้ค่า \Rsq\ ประมาณ 0 บนชุดทดสอบ และอาจติดลบเล็กน้อยหากค่าเฉลี่ยของชุดฝึกต่างจากชุดทดสอบ

\subsection{Random Forest Regressor}

Random Forest เป็นแบบจำลองรวมกลุ่ม (Ensemble Learning) ชนิด Bagging ที่พัฒนาโดย Breiman (2001) หลักการคือการสร้างต้นไม้ตัดสินใจจำนวนมากจากตัวอย่างที่สุ่มเลือกแบบมีการแทนที่ (Bootstrap Sampling) และในแต่ละโหนดจะพิจารณาเพียงคุณลักษณะย่อยที่สุ่มเลือกมา แล้วนำผลทำนายของทุกต้นมาเฉลี่ยกัน

การรวมผลจากต้นไม้หลายต้นช่วยลดความแปรปรวน (Variance) และลดปัญหาการเรียนรู้เกิน (Overfitting) ได้อย่างมีประสิทธิภาพ นอกจากนี้ค่าเบี่ยงเบนมาตรฐานของการทำนายจากต้นไม้แต่ละต้นยังสามารถใช้เป็นตัวแทนความไม่แน่นอนของการทำนายได้

\subsection{XGBoost Regressor}

XGBoost (Extreme Gradient Boosting) พัฒนาโดย Chen และ Guestrin (2016) เป็นอัลกอริทึม Boosting ที่สร้างต้นไม้แบบต่อเนื่อง โดยต้นไม้แต่ละต้นจะเรียนรู้เพื่อแก้ไขความคลาดเคลื่อนตกค้าง (Residuals) ของต้นไม้ก่อนหน้า ผ่านการปรับลดด้วยวิธี Gradient Descent

XGBoost มีกลไกการปรับให้เป็นปกติ (Regularization) ทั้งแบบ L1 และ L2 ในฟังก์ชันวัตถุประสงค์ ทำให้ควบคุมความซับซ้อนของแบบจำลองได้ดี และมีการปรับแต่งเชิงวิศวกรรมที่ทำให้ประมวลผลได้รวดเร็ว จึงเป็นที่นิยมอย่างสูงในงานข้อมูลตาราง

\subsection{Support Vector Regressor (SVR)}

SVR เป็นอัลกอริทึมที่หาไฮเปอร์เพลนที่เหมาะสมที่สุดภายในขอบเขตความคลาดเคลื่อน $\varepsilon$ (Epsilon-insensitive tube) โดยจุดข้อมูลที่อยู่ภายในขอบเขตนี้จะไม่ถูกนับเป็นความคลาดเคลื่อน การใช้ฟังก์ชันเคอร์เนลแบบ RBF (Radial Basis Function) ทำให้ SVR สามารถจับความสัมพันธ์ที่ไม่เป็นเชิงเส้นได้

SVR มีความไวต่อสเกลของข้อมูลอย่างมาก เนื่องจากการคำนวณระยะทางในฟังก์ชันเคอร์เนล จึงจำเป็นต้องปรับมาตรฐานคุณลักษณะเชิงตัวเลขก่อนใช้งาน

\section{การป้องกันการรั่วไหลของข้อมูล}

\subsection{นิยามและประเภทของ Data Leakage}

Data Leakage คือการที่สารสนเทศจากชุดข้อมูลทดสอบรั่วไหลเข้าสู่กระบวนการฝึกแบบจำลอง ทำให้ผลการประเมินสูงเกินความเป็นจริง ในบริบทของงาน QSAR มี Data Leakage ที่พบบ่อยสองรูปแบบ

\textbf{รูปแบบที่ 1: Preprocessing Leakage} เกิดจากการคำนวณค่าสถิติสำหรับการปรับมาตรฐาน เช่น ค่าเฉลี่ยและส่วนเบี่ยงเบนมาตรฐาน จากข้อมูลทั้งชุดก่อนแบ่ง Train/Test ทำให้ค่าสถิติของชุดทดสอบมีอิทธิพลต่อการแปลงข้อมูลฝึก

\textbf{รูปแบบที่ 2: Feature Selection Leakage} เกิดจากการคัดเลือกคุณลักษณะโดยพิจารณาความสัมพันธ์กับตัวแปรตามบนข้อมูลทั้งชุดก่อนแบ่งข้อมูล ทำให้คุณลักษณะที่ถูกเลือกมีความรู้เกี่ยวกับชุดทดสอบติดมาด้วย รูปแบบนี้มักส่งผลรุนแรงกว่ารูปแบบแรก โดยเฉพาะเมื่อจำนวนคุณลักษณะตั้งต้นมีมาก

\subsection{การแก้ปัญหาด้วย Pipeline และ ColumnTransformer}

Scikit-learn Pipeline เป็นกลไกที่ผนวกขั้นตอนการแปลงข้อมูลเข้ากับแบบจำลองให้เป็นวัตถุเดียว เมื่อเรียกใช้กับกระบวนการ Cross-Validation ระบบจะเรียก \texttt{fit()} ของทุกขั้นตอนใหม่ในแต่ละ Fold โดยใช้เฉพาะข้อมูลชุดฝึกของ Fold นั้น จึงรับประกันได้ว่าค่าสถิติที่ใช้แปลงข้อมูลไม่ได้มาจากชุดตรวจสอบ

\texttt{ColumnTransformer} ทำหน้าที่กำหนดว่าคอลัมน์กลุ่มใดจะผ่านการแปลงแบบใด ในโครงงานนี้กำหนดให้คอลัมน์ Fingerprint จำนวน 1{,}024 คอลัมน์ผ่านแบบ \texttt{passthrough} คือไม่แปลง และคอลัมน์ Lipinski จำนวน 4 คอลัมน์ผ่าน \texttt{StandardScaler}

\section{การประเมินประสิทธิภาพแบบจำลอง}

\subsection{ตัวชี้วัด}

\begin{table}[H]
\centering
\caption{ตัวชี้วัดประสิทธิภาพที่ใช้ในโครงงาน}
\label{tab:metrics}
\begin{tabular}{L{2.4cm} C{4.0cm} L{7.0cm}}
\toprule
\textbf{ตัวชี้วัด} & \textbf{สูตร} & \textbf{การตีความ}\\
\midrule
\Rsq & $1 - \dfrac{SS_{res}}{SS_{tot}}$ & สัดส่วนความแปรปรวนที่แบบจำลองอธิบายได้ ค่าสูงสุด 1.0 ค่าติดลบหมายถึงแย่กว่าการทำนายด้วยค่าเฉลี่ย\\
RMSE & $\sqrt{\frac{1}{n}\sum(y_i - \hat{y}_i)^2}$ & รากที่สองของค่าเฉลี่ยความคลาดเคลื่อนกำลังสอง ให้น้ำหนักกับความคลาดเคลื่อนขนาดใหญ่\\
MAE & $\frac{1}{n}\sum|y_i - \hat{y}_i|$ & ค่าเฉลี่ยความคลาดเคลื่อนสัมบูรณ์ ทนต่อค่าผิดปกติมากกว่า RMSE\\
Spearman $\rho$ & สหสัมพันธ์อันดับ & ความสามารถในการจัดอันดับสาร ซึ่งสำคัญกว่าค่าสัมบูรณ์ในงานคัดกรองยา\\
\bottomrule
\end{tabular}
\end{table}

ในบริบท QSAR ค่า RMSE และ MAE มีหน่วยเป็น log unit ค่า RMSE = 0.5 หมายความว่าโดยเฉลี่ยแบบจำลองทำนายคลาดเคลื่อนประมาณ 3 เท่าของค่า \ICfifty\ จริง

\subsection{K-Fold Cross-Validation}

การแบ่งชุดข้อมูลฝึกออกเป็น $K$ ส่วนเท่า ๆ กัน แล้วฝึกแบบจำลอง $K$ ครั้ง โดยแต่ละครั้งใช้ $K-1$ ส่วนเป็นข้อมูลฝึกและ 1 ส่วนเป็นข้อมูลตรวจสอบ ผลลัพธ์ที่ได้จึงมีความน่าเชื่อถือมากกว่าการแบ่งครั้งเดียว และค่าส่วนเบี่ยงเบนมาตรฐานข้าม Fold ยังบ่งบอกความเสถียรของแบบจำลอง

ส่วนต่างระหว่าง \Rsq\ ของชุดฝึกกับชุดตรวจสอบ (Train$-$CV Gap) เป็นดัชนีบ่งชี้การเรียนรู้เกิน ค่ายิ่งสูงแสดงว่าแบบจำลองจดจำข้อมูลฝึกมากเกินไป

\subsection{Scaffold Split}

Bemis และ Murcko (1996) เสนอแนวคิดการสกัดโครงสร้างแกนกลางของโมเลกุล (Molecular Framework) โดยตัดหมู่แทนที่ด้านข้างออก เหลือเฉพาะระบบวงแหวนและสายเชื่อม

Sheridan (2013) ชี้ว่าการแบ่งข้อมูลแบบสุ่มในงาน QSAR ทำให้สารในอนุกรมเคมีเดียวกัน (Analog Series) กระจายอยู่ทั้งชุดฝึกและชุดทดสอบ ส่งผลให้ค่า \Rsq\ สูงเกินจริง เพราะแบบจำลองเพียงต้องจดจำอนุกรมนั้นแทนที่จะเรียนรู้ความสัมพันธ์เชิงโครงสร้างที่แท้จริง การแบ่งข้อมูลโดยจัดโมเลกุลที่มี Scaffold เดียวกันไว้ฝั่งเดียวกันจึงสะท้อนสถานการณ์การทำนายสารใหม่ได้ตรงกว่า

\subsection{Applicability Domain}

Applicability Domain คือขอบเขตของปริภูมิเคมีที่แบบจำลองมีความน่าเชื่อถือ วิธีประเมินที่นิยมคือการคำนวณความคล้ายคลึงแบบ Tanimoto ระหว่างสารที่ต้องการทำนายกับสารที่ใกล้ที่สุดในชุดข้อมูลฝึก ดังสมการที่ \ref{eq:tanimoto}

\begin{equation}
T(A,B) = \frac{|A \cap B|}{|A \cup B|}
\label{eq:tanimoto}
\end{equation}

โดยทั่วไปหากค่าความคล้ายคลึงสูงสุดต่ำกว่า 0.4 ถือว่าสารนั้นอยู่นอกขอบเขตการใช้งาน และผลทำนายไม่ควรนำไปใช้เป็นค่าสัมบูรณ์

\section{การอธิบายแบบจำลองด้วย SHAP}

SHAP (SHapley Additive exPlanations) พัฒนาโดย Lundberg และ Lee (2017) เป็นวิธีอธิบายการตัดสินใจของแบบจำลองบนพื้นฐานทฤษฎีเกมสหกรณ์ โดยคำนวณค่า Shapley ซึ่งเป็นการกระจายผลตอบแทนให้ผู้เล่นแต่ละคนอย่างยุติธรรมตามการมีส่วนร่วม

ในบริบทของการเรียนรู้ของเครื่อง ผู้เล่นคือคุณลักษณะแต่ละตัว และผลตอบแทนคือค่าที่แบบจำลองทำนาย ค่า SHAP ของคุณลักษณะ $j$ ในตัวอย่าง $i$ เขียนแทนด้วย $\phi_j^{(i)}$ ซึ่งบอกว่าคุณลักษณะนั้นผลักดันการทำนายให้สูงขึ้นหรือต่ำลงเท่าใดเมื่อเทียบกับค่าฐาน ความสำคัญโดยรวมของคุณลักษณะคำนวณจากค่าเฉลี่ยสัมบูรณ์ดังสมการที่ \ref{eq:shap}

\begin{equation}
\mathrm{Importance}_j = \frac{1}{N}\sum_{i=1}^{N}\left|\phi_j^{(i)}\right|
\label{eq:shap}
\end{equation}

สำหรับแบบจำลองประเภทต้นไม้ สามารถใช้ \texttt{TreeExplainer} ซึ่งคำนวณค่า SHAP ได้อย่างแม่นยำในเวลาพหุนาม แทนที่จะต้องประมาณค่า

Rodríguez-Pérez และ Bajorath (2020) ได้ประยุกต์ใช้ค่า Shapley ในการตีความการทำนายฤทธิ์ของสารประกอบ และแสดงให้เห็นว่าวิธีนี้สามารถเชื่อมโยงผลการทำนายกลับไปยังโครงสร้างย่อยที่ตีความเชิงเคมีได้

\section{งานวิจัยที่เกี่ยวข้อง}

Nantasenamat และคณะ (2010) ได้สรุปภาพรวมเชิงปฏิบัติของงาน QSAR โดยเน้นความสำคัญของการเตรียมข้อมูล การเลือกคุณลักษณะ และการตรวจสอบความถูกต้องของแบบจำลอง งานนี้เป็นกรอบแนวคิดพื้นฐานที่โครงงานปัจจุบันยึดถือ

Gaulton และคณะ (2017) ได้อธิบายโครงสร้างและการใช้งานฐานข้อมูล ChEMBL ซึ่งเป็นฐานข้อมูลสารออกฤทธิ์ทางชีวภาพที่ผ่านการดูแลด้วยมือ (Manually Curated) ครอบคลุมข้อมูลจากวรรณกรรมเคมียาและสิทธิบัตร ประเด็นสำคัญที่งานนี้ชี้ไว้คือ ข้อมูลใน ChEMBL มาจากหลายห้องปฏิบัติการที่ใช้โปรโตคอลต่างกัน จึงมีความแปรปรวนภายในสูง จำเป็นต้องกรองและยุบข้อมูลซ้ำอย่างระมัดระวัง

Maggiora (2006) ได้อธิบายปรากฏการณ์ Activity Cliff ซึ่งเป็นกรณีที่โมเลกุลมีโครงสร้างคล้ายกันมากแต่มีฤทธิ์แตกต่างกันอย่างมาก ปรากฏการณ์นี้เป็นข้อจำกัดเชิงหลักการของแบบจำลอง QSAR ที่ตั้งอยู่บนสมมติฐานว่าโครงสร้างที่คล้ายกันควรมีฤทธิ์ใกล้เคียงกัน

\begin{keynote}[ช่องว่างที่โครงงานนี้เติมเต็ม]
งานวิจัย QSAR จำนวนมากรายงานเฉพาะผลการแบ่งข้อมูลแบบสุ่มและไม่ได้ประเมินขอบเขตการใช้งานของแบบจำลอง โครงงานนี้จึงรายงานผลทั้งการแบ่งแบบสุ่มและแบบ Scaffold ควบคู่กัน พร้อมทดสอบข้ามชุดข้อมูลกับสารจากแหล่งอิสระ และวิเคราะห์ Applicability Domain เชิงปริมาณ เพื่อให้ได้ภาพความสามารถทำนายที่ตรงกับความเป็นจริงในทางปฏิบัติ
\end{keynote}
